\section*{Sets and Indices}

\begin{itemize}
    \item $I=\{1,\dots,6\}$: set of container types, indexed by $i$ and $j$.
    \item For each demand type $j \in I$, the admissible substitute production types are
    \[
    A(j)=\{\, i \in I \mid v_i \ge v_j \,\},
    \]
    i.e., container type $i$ can satisfy demand of type $j$ only if its volume is at least as large.
\end{itemize}

\section*{Parameters}

\begin{itemize}
    \item $v_i$: volume (cm$^3$) of container type $i$.
    \item $d_i$: market demand (units) for container type $i$.
    \item $c_i$: unit variable production cost (Yuan/unit) for container type $i$.
    \item $F$: fixed setup cost (Yuan) incurred if production equipment for a type is activated.
    \item $M$: big-$M$ constant used in setup linking constraints, defined in the code as
    \[
    M=\sum_{k \in I} d_k.
    \]
\end{itemize}

Using the provided data, $F=1200$ and
\[
(v_i,d_i,c_i)_{i \in I}=
\{(1500,500,5),(2500,550,8),(4000,700,10),(6000,900,12),(9000,400,16),(12000,300,18)\}.
\]

\section*{Decision Variables}

\begin{itemize}
    \item $x_{ij}$ (code: \texttt{x[i,j]}): nonnegative integer number of units of container type $i$ produced to fulfill demand of type $j$, defined only for pairs $(i,j)$ such that $v_i \ge v_j$.
    \[
    x_{ij} \in \mathbb{Z}_{\ge 0} \qquad \forall i \in I,\ \forall j \in I \text{ with } v_i \ge v_j.
    \]
    \item $y_i$ (code: \texttt{y[i]}): binary setup variable equal to 1 if container type $i$ is produced, 0 otherwise.
    \[
    y_i \in \{0,1\} \qquad \forall i \in I.
    \]
\end{itemize}

\section*{Objective}

Minimize total variable production cost plus total fixed setup cost:
\[
\min \;
\sum_{j \in I}\sum_{i \in A(j)} c_i\, x_{ij}
\;+\;
\sum_{i \in I} F\, y_i.
\]

\section*{Constraints}

\begin{enumerate}
    \item \textbf{Demand satisfaction.} Each demand type must be met by its own type or a larger-volume substitute:
    \[
    \sum_{i \in A(j)} x_{ij} \ge d_j
    \qquad \forall j \in I.
    \]

    \item \textbf{Setup-production linking.} Production of type $i$ is allowed only if its setup is activated:
    \[
    \sum_{\substack{j \in I:\\ v_i \ge v_j}} x_{ij} \le M\, y_i
    \qquad \forall i \in I.
    \]

    \item \textbf{Integrality and domains.}
    \[
    x_{ij} \in \mathbb{Z}_{\ge 0}
    \qquad \forall (i,j): v_i \ge v_j,
    \]
    \[
    y_i \in \{0,1\}
    \qquad \forall i \in I.
    \]
\end{enumerate}

\section*{Notes on Implementation}

\begin{itemize}
    \item The implementation creates variables only for feasible substitution pairs $(i,j)$ with $v_i \ge v_j$; infeasible pairs do not exist in the model.
    \item Demand constraints are modeled with ``$\ge$'' rather than equality. Thus the implemented model permits over-satisfying a demand type, although positive costs make unnecessary overproduction unattractive in optimal solutions.
    \item There are no capacity constraints, raw material constraints, or upper bounds on total production other than the setup link through $M$.
    \item The big-$M$ value is set exactly to the sum of all demands, $M=\sum_{k\in I} d_k$, matching the code.
    \item The model is solved as a mixed-integer linear program using the Gurobi solver through PuLP-compatible calls.
\end{itemize}
